\chapter{Experiments}

%%%%%%%%%%%%%%%%%%%%%%%%%%%%%%%%%%%%%%%%%%%%%%%%%
\section{Monomial Approximation}

\subsubsection{Support in $[0,1]$ and $N=M$}
We have evaluated the two techniques described in Section~\ref{subsec:monomial_approximation_support_0_1_N_equal_M}. For both methods the tests include the same polynomial expressed in monomial form, i.e. $\displaystyle p_n(x)=\sum_{i=0}^{n} a_i x^i$, and different orders $n$. These experiments are in \href{https://github.com/edoardosarri24/probabilistic-soft-real-time-analysis/blob/main/bernstein/src/main/java/Main_monomial_approximation_support_0_1_N_equal_M.java}{this Main.java}.

The results are aligned with the limitations of the \texttt{double} type described in Section~\ref{sec:bernstein_library_constraints}.

\myskip

If we look at the first method, i.e. the direct transformation of the monomial coefficients to Bernstein coefficients, the results confirm that the method is exact; the two polynomials, the approximator and the approximatee are the same. We can see that the error is in the order of $10^{-16}$: if degree is 25 the $L_1$ norm error is $2\times10^{-15}$; if the degree is 1029 the error is $L_1=9.7\times10^{-15}$. If we increase further the degree by one unit, the system fails and we obtain a $NaN$, exactly as described in Section~\ref{sec:bernstein_library_constraints}.

\myskip

If we look to the method that use the basis matrix to compute the Bernstein coefficients, the things change drastically. Already for a small degree, like 10, we have an error of $L_1=1.84 \times 10^{-14}$; if we increase the degree to 35 we have $L_1=2.04\times 10^{-2}$; if the degree is greater than 50 the problem becomes intractable and we obtain an error of $L_1=1.7\times10^{5}$.